\documentclass[11pt]{article}

\usepackage[margin=1in]{geometry}
\usepackage{amsmath}
\usepackage{amssymb}
\usepackage{booktabs}
\usepackage{siunitx}
\usepackage{microtype}
\usepackage[hidelinks]{hyperref}

\emergencystretch=3em

\newcommand{\ii}{\mathrm{i}}
\newcommand{\dd}{\,\mathrm{d}}
\newcommand{\Rw}{R_{\mathrm{w}}}
\DeclareMathOperator{\sech}{sech}

\title{Beyond the Symmetric Thin Ship:\\
Asymmetry, Centreplane Lifting, and Heel\\
\large (Part~II of the \texttt{michell} calculation notes)}
\author{\texttt{michell}}
\date{}

\begin{document}
\maketitle

\begin{abstract}
Part~I documented the classical Michell wave-resistance integral for a hull
that is symmetric about its centerplane: a thin \emph{thickness} distribution
radiating a \emph{source} wave system, evaluated exactly per B-spline knot span.
This second note documents the three extensions built on top of that kernel, all
of which reuse the Part~I machinery essentially unchanged. First, a
\emph{port/starboard-asymmetric} hull is decomposed into a symmetric thickness
part and an antisymmetric \emph{camber} part; the camber radiates a centreplane
\emph{y-dipole} wave system that superposes on the source system without a cross
term. Second, the dipole \emph{magnitude} --- which no pointwise boundary
condition fixes, because the doublet density is a non-local functional of the
camber slope --- is obtained honestly by \emph{solving} a centreplane
lifting-surface problem with a horseshoe vortex lattice, replacing the crude
strip closure $\mu = 2U f_a$. Third, a \emph{heeled} hull is handled by tilting
the source centreplane, which turns the vertical decay rate \emph{complex} and
otherwise leaves the exact per-span kernel untouched. Throughout, section
references point at the Rust source. All notation and constants are inherited
from Part~I; in particular $\nu = g/U^2$, $k_x = \nu\lambda$,
$\kappa = \nu\lambda^2$, and $F(\lambda) = I + \ii J$ is the source free-wave
amplitude of Part~I, Eq.~(inner).
\end{abstract}

\tableofcontents

\section{Why the symmetric integral is not enough}
\label{sec:motivation}

The classical Michell integral of Part~I represents the hull by a single-valued
half-breadth $f(x,z)\ge 0$ and a sheet of \emph{sources} of strength
$\sigma = 2U\,\partial f/\partial x$ on the centerplane. That model is blind to
three physically important situations:

\begin{enumerate}
  \item \textbf{Port/starboard asymmetry.} A hull whose two sides differ ---
        a cambered demihull, an asymmetric catamaran float, a hull at a small
        drift (yaw) angle --- carries a transverse \emph{lift} in addition to
        its thickness. Its wave system has an antisymmetric component that a
        single half-breadth cannot express.
  \item \textbf{The lifting closure.} The strength of the antisymmetric
        (dipole) system is \emph{not} fixed pointwise by the body boundary
        condition, unlike the source strength. Getting its magnitude right
        requires solving a lifting-surface integral equation.
  \item \textbf{Heel.} A hull rolled about its longitudinal axis is asymmetric
        \emph{relative to the horizontal free surface}, even if it is
        geometrically symmetric about its own centreplane. Its keel swings off
        the earth-vertical plane and it can no longer be written as
        port/starboard half-beams there.
\end{enumerate}

The unifying idea is that all three are small perturbations that enter through
the \emph{same} inner integral and the \emph{same} outer $\theta$-quadrature as
Part~I. The only new ingredients are (a) a second coefficient field for the
antisymmetric slope, (b) a real spectral \emph{weight} coupling that field to a
transverse-wavenumber factor, (c) a linear solve for the lifting density, and
(d) a complex generalisation of the vertical-decay moment. Table~\ref{tab:map}
maps each extension to its entry point in the crate.

\begin{table}[h]
  \centering
  \begin{tabular}{lll}
    \toprule
    Extension & Entry point & Kernel change \\
    \midrule
    Asymmetry (strip closure)
      & \texttt{multihull\_wave\_resistance\_with}
      & second slope field $\partial f_a/\partial x$ \\
    Asymmetry (solved lift)
      & \texttt{asymmetric\_wave\_resistance\_lifting}
      & centreplane vortex lattice \\
      & \texttt{multihull\_wave\_resistance\_lifting} & \\
    Heel
      & \texttt{heel\_wave\_resistance}
      & complex vertical decay $\kappa$ \\
    \bottomrule
  \end{tabular}
  \caption{The three extensions and where they attach to the Part~I kernel
  (\texttt{michell.rs}).}
  \label{tab:map}
\end{table}

\section{Decomposition of an asymmetric hull}
\label{sec:decomp}

An asymmetric hull is built from \emph{two} half-breadth surfaces
(\texttt{Hull::new\_asymmetric}, \texttt{hull.rs}): a starboard side
$y = +f_+(x,z)\ge 0$ and a port side $y = -f_-(x,z)\le 0$, both stored as
non-negative magnitudes. The two surfaces must share identical degrees and knot
vectors --- only their control nets may differ --- so that both live on the same
grid of knot spans. On that shared parametrisation the hull splits linearly into
a \emph{symmetric} (thickness) part and an \emph{antisymmetric} (camber) part,
\begin{equation}
  f_{\mathrm{sym}}(x,z) = \tfrac12\bigl(f_+ + f_-\bigr),
  \qquad
  f_a(x,z) = \tfrac12\bigl(f_+ - f_-\bigr).
  \label{eq:split}
\end{equation}
Because a B-spline is linear in its control net, \eqref{eq:split} is applied
directly to the control points: $P^{\mathrm{sym}}_{ij} = \tfrac12(P^{+}_{ij} +
P^{-}_{ij})$ and $P^{a}_{ij} = \tfrac12(P^{+}_{ij} - P^{-}_{ij})$. The symmetric
net is itself a valid non-negative half-breadth (a mean of two non-negative
nets), so the base hull is constructed from it and \emph{every} validated
symmetric quantity of Part~I --- volume, LCB/VCB, the source slope field
$c_{ab}$ --- is reused verbatim. The antisymmetric net produces a second
per-span slope field
\begin{equation}
  \frac{\partial f_a}{\partial x}(x,z)
    = \sum_{a,b} c^{a}_{ab}\,(x-x_0)^a (z-z_0)^b,
  \qquad
  \text{(\texttt{fx\_a\_coeff}),}
  \label{eq:fxa}
\end{equation}
computed by exactly the same \texttt{compute\_fx\_coeff} routine as the source
field, and stored with the identical flattened layout. When $f_+ \equiv f_-$ the
antisymmetric net is zero and the hull is bit-for-bit the symmetric
\texttt{Hull::new} of Part~I.

Two integrated quantities do \emph{not} reduce to twice the mean and are
corrected for the two sides separately: the wetted surface area and the
centreplane transverse (waterplane) inertia both see $f_+$ and $f_-$
individually. Volume, LCB/VCB, and waterplane area depend only on the sum
$f_+ + f_- = 2 f_{\mathrm{sym}}$ and are already correct from the base hull.

\section{The centreplane dipole wave system}
\label{sec:dipole}

\subsection{Source plus dipole}

The symmetric part $f_{\mathrm{sym}}$ radiates the classical source system with
the Part~I free-wave amplitude
\begin{equation}
  F(\lambda) = \iint \frac{\partial f_{\mathrm{sym}}}{\partial x}\,
    e^{-\kappa z}\, e^{\,\ii k_x x}\dd x\dd z,
  \qquad k_x = \nu\lambda,\; \kappa = \nu\lambda^2.
\end{equation}
The antisymmetric part $f_a$ is a \emph{camber} distribution on the centreplane.
Physically it is a sheet of transverse (\,$y$-normal\,) \emph{doublets} of density
$\mu(x,z)$: the potential jump across the centreplane. Such a doublet sheet
radiates a free wave carrying \emph{one extra factor of the transverse
wavenumber} $k_y = \nu\lambda\sqrt{\lambda^2-1}$ relative to a source of the same
integrated strength,
\begin{equation}
  A_d(\theta) = \ii\,k_y \iint \mu(x,z)\, e^{-\kappa z}\, e^{\,\ii k_x x}
    \dd x\dd z.
  \label{eq:Ad-raw}
\end{equation}
Everything downstream is a matter of fixing $\mu$ and reducing
\eqref{eq:Ad-raw} to a weight on a familiar integral.

\subsection{The strip closure and the real dipole weight}
\label{sec:stripweight}

The crudest way to fix the density is the \emph{strip closure} $\mu = 2U f_a$
--- the direct antisymmetric analogue of Michell's source strength
$\sigma = 2U\,\partial f_{\mathrm{sym}}/\partial x$. Substituting it into
\eqref{eq:Ad-raw} and integrating by parts in $x$, using that $f_a$ closes
(vanishes) at bow and stern so the boundary term drops,
\begin{equation}
  \iint f_a\, e^{-\kappa z} e^{\,\ii k_x x}\dd x\dd z
   = \frac{\ii}{\nu\lambda}
     \underbrace{\iint \frac{\partial f_a}{\partial x}\,
       e^{-\kappa z} e^{\,\ii k_x x}\dd x\dd z}_{\textstyle G(\lambda)},
  \label{eq:ibp}
\end{equation}
which introduces the companion amplitude $G(\lambda)$ --- the \emph{same} inner
integral as the source, but over the antisymmetric slope field \eqref{eq:fxa}.
The two factors of $\ii$ (one from the dipole, one from the integration by
parts) collapse to a \emph{real} weight in this crate's amplitude units
(dividing by $2U$),
\begin{equation}
  \boxed{\;
  A_d(\lambda) = -\,w(\lambda)\,G(\lambda),
  \qquad
  w(\lambda) = C_{\mathrm{dip}}\,\sqrt{\lambda^2 - 1}. \;}
  \label{eq:dipoleweight}
\end{equation}
This is \texttt{dipole\_weight} in the source; the closure constant
$C_{\mathrm{dip}} = 1$ (\texttt{DIPOLE\_WEIGHT\_C}) is the value self-consistent
with the source normalisation under the strip closure. The weight vanishes as
$\lambda\to 1$ ($\theta\to 0$): a wave running dead ahead carries no transverse
wavenumber and is blind to side-to-side asymmetry, so the whole dipole
contribution is fed by the diverging (large-$\lambda$) part of the spectrum.

\paragraph{Warning on magnitude.} The \emph{structure} of
\eqref{eq:dipoleweight} --- the $\sqrt{\lambda^2-1}$ weight, the odd-in-$\theta$
parity, the additive separation below --- is exact. The overall
\emph{constant} $C_{\mathrm{dip}}$ is not. Unlike the source strength, which the
boundary condition fixes pointwise (a source sheet's normal-velocity
\emph{jump} equals its local strength), the doublet density is fixed by the
\emph{mean} of the two-sided normal velocities, and a doublet sheet's mean
normal velocity is a hypersingular (finite-part) integral of $\mu$. So $\mu$ is
genuinely \emph{non-local} in $\partial f_a/\partial x$: no exact pointwise
closure exists, and $\mu = 2U f_a$ is only the crudest strip estimate. Getting
the magnitude right is the subject of Section~\ref{sec:lifting}.

\subsection{Additive resistance and the $\theta$-parity argument}
\label{sec:parity}

Combining the two systems, hull member $j$ (at fleet offset $\Delta x_j$,
transverse position $y_j$) contributes to the $\pm\theta$ wave systems the
amplitudes
\begin{equation}
  A_{\pm,j} = \bigl(F_j \mp w\,G_j\bigr)\,
    \exp\!\bigl[\ii\,(k_x\Delta x_j \pm k_y y_j)\bigr],
\end{equation}
and the resistance integrand is the mean of the two signed systems, exactly as
for a symmetric fleet in Part~I:
\begin{equation}
  \Rw = \frac{4\rho g^2}{\pi U^2}\int_0^{\pi/2}
    \tfrac12\bigl(|A_+|^2 + |A_-|^2\bigr)\,\sec^3\theta\dd\theta,
  \label{eq:Rasym}
\end{equation}
with $A_\pm = \sum_j A_{\pm,j}$ (\texttt{amp\_sq} in
\texttt{multihull\_wave\_resistance\_with}). The source amplitude $F$ is
\emph{even} in $\theta$ and the dipole term $wG$ is \emph{odd} (the $+\theta$
system carries $-wG$, the $-\theta$ system $+wG$). Their cross term therefore
integrates to zero over the symmetric Kelvin fan, and for a single hull
\eqref{eq:Rasym} separates cleanly:
\begin{equation}
  \boxed{\;\Rw = R_{\mathrm{source}}(f_{\mathrm{sym}})
    + R_{\mathrm{dipole}}(f_a).\;}
  \label{eq:additive}
\end{equation}
A symmetric hull ($f_a\equiv 0$) recovers classical Michell exactly. For a fleet
the dipole amplitudes superpose and interfere with the source amplitudes through
the placement phases, so an asymmetric-demihull catamaran's asymmetry and its
demihull interference are handled together in one spectrum.

\paragraph{Kernel cost.} The source and dipole amplitudes share a single pass of
the $z$-moments (Part~I, Eqs.~(N)): only the $x$-span accumulation is repeated
against the second coefficient array. \texttt{InnerIntegral::eval\_pair} returns
$(F, G)$ together, so an asymmetric evaluation costs barely more than a
symmetric one.

\section{The centreplane lifting solve (``level~2'')}
\label{sec:lifting}

To replace the strip constant with a physically grounded density we
\emph{solve} the lifting-surface problem for $\mu(x,z)$. The centreplane
$x\in[0,L]$, $z\in[0,T]$ carries the camber surface $f_a$, whose streamwise
slope $\partial f_a/\partial x$ is the downwash the transverse doublet sheet must
produce --- the direct analogue of a wing's camber-slope condition.

\subsection{Two-dimensional proving ground}
\label{sec:2d}

Before the delicate 3D free-surface solve, the crate builds and validates its
2D ancestor --- classical thin-airfoil theory (\texttt{lifting.rs}) --- which
exercises the same machinery: discretising a lifting sheet, imposing the Kutta
condition, and recovering the loading. A thin airfoil of unit chord with
camber-line slope $\dd\eta/\dd x$ is represented by a bound-vortex sheet
$\gamma(x)$ satisfying the Cauchy singular integral equation
\begin{equation}
  \frac{1}{2\pi}\;\text{p.v.}\!\!\int_0^1
    \frac{\gamma(\xi)}{x - \xi}\dd\xi
    = U\Bigl(\alpha - \frac{\dd\eta}{\dd x}(x)\Bigr).
\end{equation}
Two independent solvers cross-check each other and closed-form results:
\begin{itemize}
  \item \textbf{Glauert Fourier series} (\texttt{glauert}), with
        $x = \tfrac12(1-\cos\theta)$ and
        $\gamma/U = 2\bigl[A_0\frac{1+\cos\theta}{\sin\theta} +
        \sum_{k\ge1}A_k\sin k\theta\bigr]$; the $\sin k\theta$ terms vanish at
        both edges so the Kutta condition is built in, and
        $C_L = \pi(2A_0 + A_1)$. A flat plate gives $A_0 = \alpha$, hence the
        canonical $C_L = 2\pi\alpha$.
  \item \textbf{Lumped-vortex lattice} (\texttt{solve\_vortex\_lattice}) with
        the Pistolesi $\tfrac14$--$\tfrac34$ rule: a point vortex at each
        panel's quarter-point, flow tangency collocated at its three-quarter
        point, which enforces the Kutta condition automatically. This is the
        desingularised 2D form of the hypersingular kernel the 3D solve must
        integrate.
\end{itemize}
In hull terms the camber line is the antisymmetric half-beam $f_a$ at a
waterline strip and the forcing is $-U\,\partial f_a/\partial x$ (no separate
incidence). The solved loading integrates to the doublet density
$\mu(x) = \int_0^x \gamma\dd\xi$ --- exactly the quantity the free-wave
amplitude carries.

\subsection{Three-dimensional horseshoe vortex lattice}
\label{sec:3d}

The 3D engine (\texttt{lifting3d.rs}) generalises the lattice: each panel carries
a \emph{horseshoe vortex} --- a bound segment at the quarter-chord plus two
trailing legs running downstream to $+\infty$ along $\hat{x}$ --- and flow
tangency is collocated at the three-quarter-chord point. The trailing legs make
the Kutta condition automatic. Assembling the normal-wash influence
coefficients $A_{ij} = (\mathbf{v}_{j\to i}\cdot\hat{n}_i)$ from the
Biot--Savart law (finite-segment and semi-infinite closed forms, after Katz \&
Plotkin) and solving
\begin{equation}
  A\,\Gamma = b, \qquad b_i = -\,\mathbf{U}_\infty\cdot\hat{n}_i,
\end{equation}
gives the bound circulation $\Gamma$ on every panel. Lift follows from
Kutta--Joukowski, $C_L = 2\sum_j \Gamma_j\,\Delta y_j / S$. The engine is
validated on a flat rectangular wing against the two exact asymptotic limits
that bracket the lift-curve slope: $\dd C_L/\dd\alpha\to 2\pi$ as
$\mathrm{AR}\to\infty$ (2D limit) and $\to \tfrac{\pi}{2}\mathrm{AR}$ as
$\mathrm{AR}\to 0$ (R.T.\ Jones slender-wing limit), with the Prandtl
lifting-line estimate $2\pi/(1+2/\mathrm{AR})$ in between.

\subsection{Free surface as a rigid wall (double body)}
\label{sec:doublebody}

The hull-specific step (\texttt{centerplane.rs}) applies this lattice to the
\emph{vertical} centreplane. In the low-Froude near field the free surface
$z=0$ behaves as a rigid wall ($\partial\phi/\partial z = 0$), consistent with
how Michell's source strength is set on the double-body flow. This is enforced
by \textbf{reflecting the centreplane about $z=0$} with the same circulation
sign, so the bound vortices run continuously through the waterline. The
consequences:
\begin{itemize}
  \item the waterline acts as a wing \emph{root} (maximum loading) and the keel
        $z=T$ as a free \emph{tip} (loading relieved to zero);
  \item a surface-piercing lifting surface of draft $T$ therefore behaves as a
        wing of span $2T$ --- the classical \emph{effective-aspect-ratio
        doubling}, checked in the tests against a full wing of span $2T$ to
        $\sim0.1\%$ at $\alpha = \SI{3}{\degree}$.
\end{itemize}
The camber enters through the panel \emph{normal}
$\hat{n} = (-f_{a,x}, 1, 0)/\sqrt{1 + f_{a,x}^2}$, so the freestream stays
$(1,0,0)$ and the solve is called at zero incidence. Solving for $\Gamma$ and
accumulating from the leading edge gives the doublet density
$\mu(x,z) = \int\gamma$ on the physical half $z\ge 0$, together with a
side-force coefficient $C_Y = 2\sum_{\mathrm{phys}}\Gamma\,\Delta z/(LT)$.

\subsection{From the solved density to the wave amplitude}
\label{sec:solvedamp}

The solved density feeds the dipole free-wave amplitude by midpoint quadrature
over the panel grid (\texttt{doublet\_free\_wave\_amplitude}, each sample
carrying cell area $\Delta x\,\Delta z$),
\begin{equation}
  G_\mu(\lambda) = \iint \frac{\mu(x,z)}{U}\,
    e^{-\kappa z}\, e^{\,\ii k_x x}\dd x\dd z,
\end{equation}
from which the dipole amplitude is
\begin{equation}
  \boxed{\;
  A_d(\lambda) = \ii\,\tfrac12\,k_y\,G_\mu(\lambda),
  \qquad k_y = \nu\lambda\sqrt{\lambda^2-1}. \;}
  \label{eq:Ad-solved}
\end{equation}
Note that here $G_\mu$ integrates $\mu$ \emph{directly} (not the slope), so no
integration by parts is applied and the amplitude stays imaginary --- contrast
the strip form \eqref{eq:dipoleweight}, where the by-parts step made it real.
The two paths share \emph{one} normalisation: substituting the strip closure
$\mu/U = 2 f_a$ into \eqref{eq:Ad-solved} and integrating by parts via
\eqref{eq:ibp} reproduces $A_d = -\sqrt{\lambda^2-1}\,G$ exactly. The factor
$\tfrac12$ is precisely what makes the closure $\mu/U = 2 f_a$ collapse onto the
$C_{\mathrm{dip}} = 1$ weight. The solved $\mu$ has a different
chordwise/vertical \emph{shape} than $2 f_a$, giving a dipole of the same order
but a speed-dependent ratio; for a 2D flat plate $\mu_{\mathrm{lift}} /
\mu_{\mathrm{strip}} = \pi/2$.

The single-hull entry point \texttt{asymmetric\_wave\_resistance\_lifting} and
the fleet form \texttt{multihull\_wave\_resistance\_lifting} solve each
asymmetric member's centreplane \emph{once} (up front, not per $\lambda$), then
superpose the source amplitude $F_j$ and the solved dipole $A_{d,j}$ with the
member's placement phase, exactly as in \eqref{eq:Rasym}. The dipole phase is
re-centred to the hull midpoint (the solver works on $x\in[0,L]$) so it shares
the source's origin. This is the far-field coupling: each member's lifting
problem is solved independently with its own free-surface image, and their
dipole \emph{wave} systems interfere through the superposition --- so an
asymmetric-demihull catamaran's demihull interference is captured. Near-field
lifting cross-induction between close demihulls is \emph{not} (that would need
one coupled lifting solve over all centreplanes at once). The default grid is
$n_x = 48$ streamwise $\times\ n_z = 16$ vertical panels (\texttt{LiftingGrid}).

\section{Heeled hulls}
\label{sec:heel}

A hull heeled by an angle $\varphi$ about its longitudinal axis is asymmetric
relative to the horizontal free surface, but its keel swings off the
earth-vertical centreplane, so it cannot be written as port/starboard half-beams
there. Thin-ship theory instead keeps the sources on the ship's \emph{own}
(now tilted) centreplane: a strip at ship-depth $z$ sits at earth depth
$z\cos\varphi$ and transverse offset $-z\sin\varphi$. The transverse offset
contributes a transverse-wavenumber phase $e^{-\ii k_y\,(-z\sin\varphi)}$ to the
kernel, which combines with the earth-vertical decay $e^{-\nu\lambda^2
z\cos\varphi}$ into a single \textbf{complex vertical decay}
\begin{equation}
  \boxed{\;
  \kappa = \nu\lambda^2\cos\varphi
    \;+\; \ii\,\nu\lambda\sqrt{\lambda^2-1}\,\sin\varphi
    \;=\; \nu\lambda^2\cos\varphi + \ii\,k_y\sin\varphi. \;}
  \label{eq:heelkappa}
\end{equation}
The imaginary part is the same transverse-wavenumber coupling as the asymmetric
dipole, arising here purely from geometry.

This is the \emph{only} change to the Part~I kernel. The $x$-oscillation moments
$M_a$ and the per-span accumulation are identical; only the $z$-moment
$N_b = \int_0^{h_z} Z^b e^{-\kappa Z}\dd Z$ becomes complex, evaluated by
\texttt{exp\_moments\_complex} --- the \emph{same} series and integration-by-parts
recurrence as the real routine (Part~I, Eqs.~(Nrec)--(Nseries)), promoted to
$\mathbb{C}$ with $\mathrm{Re}\,\kappa\ge 0$ so the decay never grows. The real
path is left entirely untouched, so $\varphi = 0$ reproduces
\texttt{wave\_resistance} exactly (\texttt{HeelInner} in \texttt{michell.rs}).

The tilt makes the amplitude $F$ depend on the \emph{sign} of $\theta$ (through
the $\pm$ in $k_y\sin\varphi$), so the resistance averages the two signed
systems,
\begin{equation}
  \Rw = \frac{4\rho g^2}{\pi U^2}\int_0^{\pi/2}
    \tfrac12\bigl(|F_+|^2 + |F_-|^2\bigr)\,\sec^3\theta\dd\theta,
\end{equation}
where $F_\pm$ uses $\kappa$ with $\pm\ii k_y\sin\varphi$. The result is
\emph{even} in $\varphi$: port and starboard heel are mirror images. Because the
tilt spreads the centreplane transversely by up to $T|\sin\varphi|$, the
outer-quadrature panels are sized with an effective transverse half-extent
$y_{1/2} = T|\sin\varphi|$ to resolve the heel-induced oscillation.

This captures the asymmetric \emph{wave-making} of the tilted thickness
distribution --- the leading heel effect. It does \emph{not} include the lifting
side-force that a heeled-and-yawed (drifting) hull develops; that would enter as
a separate forcing into the centreplane lifting solve of
Section~\ref{sec:lifting}.

\section{Data-flow summary}
\label{sec:dataflow2}

\begin{center}
\begin{tabular}{r c l}
  \toprule
  \multicolumn{3}{c}{\textbf{Asymmetry --- strip closure}} \\
  \midrule
  two nets $P^\pm_{ij}$ & $\longrightarrow$ & $f_{\mathrm{sym}}, f_a$ (Eq.~\ref{eq:split}); slope fields $c_{ab}, c^a_{ab}$ \\
  $c_{ab}, c^a_{ab}$ & $\longrightarrow$ & $(F, G)$ together (\texttt{eval\_pair}) \\
  $G$, weight $w=\sqrt{\lambda^2-1}$ & $\longrightarrow$ & dipole $A_d = -wG$ (Eq.~\ref{eq:dipoleweight}) \\
  $\tfrac12(|F{-}wG|^2 + |F{+}wG|^2)$ & $\longrightarrow$ & outer $\theta$-quadrature \;$\Rightarrow\; \Rw = R_{\mathrm{src}} + R_{\mathrm{dip}}$ \\
  \midrule
  \multicolumn{3}{c}{\textbf{Asymmetry --- solved lift}} \\
  \midrule
  $\partial f_a/\partial x$ on centreplane & $\longrightarrow$ & vortex-lattice solve $\Gamma$ (double body) \\
  $\Gamma$ & $\longrightarrow$ & doublet density $\mu = \int\gamma$, side force $C_Y$ \\
  $\mu$ & $\longrightarrow$ & $G_\mu$; dipole $A_d = \ii\tfrac12 k_y G_\mu$ (Eq.~\ref{eq:Ad-solved}) \\
  $F_j + A_{d,j}$, placements & $\longrightarrow$ & outer $\theta$-quadrature $\Rightarrow \Rw$ \\
  \midrule
  \multicolumn{3}{c}{\textbf{Heel}} \\
  \midrule
  heel $\varphi$ & $\longrightarrow$ & complex $\kappa$ (Eq.~\ref{eq:heelkappa}) \\
  $\kappa$ & $\longrightarrow$ & complex $z$-moments $N_b$ (\texttt{exp\_moments\_complex}) \\
  $\tfrac12(|F_+|^2 + |F_-|^2)$ & $\longrightarrow$ & outer $\theta$-quadrature $\Rightarrow \Rw$ (even in $\varphi$) \\
  \bottomrule
\end{tabular}
\end{center}

\noindent All three extensions share the exact per-span inner kernel and the
single smooth outer $\theta$-quadrature of Part~I. The only approximations added
are the strip closure constant $C_{\mathrm{dip}}$ (Section~\ref{sec:stripweight},
removed by the lifting solve), the panel discretisation of the lifting solve and
its midpoint quadrature (Section~\ref{sec:lifting}), and --- as in Part~I --- the
one-dimensional angular quadrature. The wave-pattern (wake) branch of Part~I is
unchanged: the free-wave elevation grid is still built from the symmetric source
amplitudes.

\paragraph{References.} Tuck (1989); Dambrine, Pierre \& Rousseaux (2016);
Tuck, Scullen \& Lazauskas (2001, 2002); Katz \& Plotkin, \emph{Low-Speed
Aerodynamics} (horseshoe-lattice \texttt{VORTXL}); Glauert, \emph{The Elements of
Aerofoil and Airscrew Theory}; Kaklis \& Papanikolaou (21st Symp.\ Naval Hydro.,
Appendix~A, on the hypersingular doublet closure).

\end{document}
